% Source: https://artofproblemsolving.com/community/c3h3176265_2019_mbmt_geometry_round__montgomery_blair_math_tournament
% Stable identity: aops-topic-3176265-q11
Put $A=(-\dfrac{1}{2},0)$, $B=(\dfrac{1}{2},0)$. Every circumcenter lies on $x=0$, so any periodic point has form $X=(0,y)$ with $y\ne0$. Equal squared distances give $f(0,y)=(0,\dfrac{y^2-\dfrac{1}{4}}{2y})$. Parameterize $y=\dfrac{1}{2}\cot\theta$ with unique $0<\theta<\pi$. The double-angle identity turns the map into $\theta\mapsto2\theta$ modulo $\pi$. Thus periodicity after $2019$ iterations requires $(2^{2019}-1)\theta=k\pi$. The possible integers are $1\le k\le2^{2019}-2$, each yielding a different point. None has an undefined intermediate iterate: a zero height would require $2^{j+1}k$ to be an odd multiple of the odd number $2^{2019}-1$, impossible by parity; a multiple of $\pi$ is also impossible because this odd denominator is coprime to every power of two. Therefore all candidates work, and the count is $\boxed{2^{2019}-2}$.
\begin{asy}
usepackage("amsmath");
size(7cm);
pair A=(-.5,0),Bp=(.5,0),X=(0,1),P=(0,.375);draw(circle(P,.625));draw(A--Bp--X--cycle);draw((0,-.4)--(0,1.3),dashed);label("$A$",A,SW);label("$B$",Bp,SE);label("$X$",X,N);dot("$f(X)$",P,E);
\end{asy}
